\section{Tổng quan nghiên cứu}
Các hướng tiếp cận cải thiện IAQ truyền thống thường dựa trên thông gió, máy lọc không khí hoặc điều khiển điều hòa. Các giải pháp này có hiệu quả nhất định nhưng có thể phát sinh chi phí năng lượng, chi phí bảo trì màng lọc và thường chưa gắn với nhu cầu tạo không gian sống xanh. Trong khi đó, các nghiên cứu về cây trồng trong nhà, thủy canh và bộ lọc thực vật cho thấy cây có thể đóng vai trò hỗ trợ trong việc giảm CO2, VOC, cải thiện độ ẩm cục bộ và tạo môi trường sinh hoạt dễ chịu hơn \cite{hydroponic_iaq_survey,montaluisa_pothos_biofilter}. Với trầu bà, tài liệu thực nghiệm về bộ lọc thực vật chủ động ghi nhận hiệu suất loại bỏ ổn định cao đối với acetone, $\alpha$-pinene và toluene trong điều kiện có tuần hoàn khí, dung dịch dinh dưỡng và hệ rễ phát triển \cite{montaluisa_pothos_biofilter}.

Các hệ thống nhà thông minh hiện nay thường tập trung vào điều khiển thiết bị như đèn, quạt, điều hòa hoặc giám sát an ninh. Một số hệ thống có bổ sung cảm biến nhiệt độ, độ ẩm và khí CO2, nhưng dữ liệu môi trường thường được xử lý tách biệt với nhu cầu chăm sóc cây trồng. Trong khi đó, các hệ thống trồng cây thông minh lại chủ yếu quan tâm đến độ ẩm đất, pH, TDS hoặc ánh sáng, chưa gắn chặt với bài toán kiểm soát không khí trong căn hộ. Khoảng tách này làm giảm giá trị của mô hình cây lọc khí, bởi sức khỏe cây và vùng rễ là một phần quan trọng của cơ chế hấp thụ, chuyển hóa và hỗ trợ vi sinh vật trong hệ lọc thực vật.

Về công nghệ, IoT cho phép kết hợp cảm biến, bộ chấp hành, vi điều khiển và nền tảng dữ liệu để giám sát môi trường theo thời gian thực \cite{wang_smart_agriculture,liu_realtime_iaq,singh_iot_air_quality}. Khi áp dụng cho trầu bà hoặc mô hình thủy canh, IoT có thể hỗ trợ điều chỉnh ánh sáng, nước, dinh dưỡng và cảnh báo bất thường, qua đó duy trì điều kiện tốt hơn cho cây và cho không gian sống. Tài liệu về quản lý pH/EC trong thủy canh nhấn mạnh rằng nồng độ muối dinh dưỡng và pH dung dịch cần được kiểm tra thường xuyên, vì sai lệch có thể gây thiếu dinh dưỡng, độc ion hoặc giảm sinh trưởng \cite{singh_hydroponics_ph_ec}. Tuy nhiên, nhiều nghiên cứu vẫn dừng ở mô hình giám sát riêng lẻ hoặc chưa nhấn mạnh luồng điều khiển hai chiều từ ứng dụng người dùng đến thiết bị tại hiện trường.

Về truyền thông, Wi-Fi là lựa chọn phổ biến cho thiết bị IoT trong nhà vì có sẵn hạ tầng mạng. Tuy nhiên, node cảm biến đặt ở ban công, gần cửa sổ hoặc khu vực nhiều vật cản có thể gặp tình trạng tín hiệu Wi-Fi không ổn định. LoRa có ưu điểm truyền xa, công suất thấp, khả năng xuyên vật cản tốt hơn trong nhiều kịch bản và phù hợp với node cảm biến gửi dữ liệu định kỳ dung lượng nhỏ. Vì vậy, khóa luận sử dụng LoRa cho liên kết node - gateway, còn gateway ESP32 đảm nhiệm kết nối Internet và đẩy dữ liệu lên ThingsBoard.

Nền tảng ThingsBoard hỗ trợ quản lý thiết bị, lưu trữ telemetry, dashboard và cơ chế RPC/attribute phục vụ điều khiển hai chiều. Ứng dụng Flutter có thể khai thác API của ThingsBoard để xây dựng giao diện di động đa nền tảng, hiển thị dữ liệu và gửi lệnh điều khiển. Cách tiếp cận này giúp tách rõ phần thu thập dữ liệu, phần truyền thông, phần nền tảng IoT và phần giao diện người dùng.

Từ các phân tích trên, khoảng trống mà khóa luận hướng đến là xây dựng một hệ thống quy mô nhỏ, chi phí hợp lý, có thể đặt trong căn hộ chung cư, đồng thời kết hợp dữ liệu IAQ với dữ liệu chăm sóc cây trầu bà. Hệ thống cần không chỉ ghi nhận thông số mà còn cho phép cảnh báo, điều khiển thủ công và làm nền tảng cho các luật điều khiển tự động trong tương lai.
